% Proposed messaging design; included by the main whitepaper.
\section{Agent-to-agent messaging across a global HF station network}
\label{sec:agent-chat}

\subsection{Scope, status, and architectural intent}
\label{sec:chat-scope}
This chapter proposes an asynchronous messaging service through which agents
at independently operated stations exchange compact text and structured
messages over eligible HF contacts. ChatGPT Agent and Codex are possible
local orchestration clients. MCP connects each agent host to its own station;
it is not the over-the-air protocol, a globally shared agent bus, or a source
of modem timing. A remote agent need not use the same model or host software.

\textbf{Implementation status.} The Rust project contains native BW500 codec,
streaming reception, framing, handshake and ARQ prototypes, radio-control MCP,
and a SigPlot audio demonstration. These do not yet constitute an integrated
operational modem or an agent messaging service. All chat interfaces, durable
stores, receipt protocols, schedulers, identity bindings, and relay behavior in
this chapter are \emph{proposed}. Native codec round trips and recorded-signal
replay are not qualification of live stock VARA, RF messaging, or global
reachability. Section~\ref{sec:modem-software} remains the implementation and
DSP evidence baseline.

The initial product objective is deliberately observable: agent A queues
``Hello'', station B stores it once, agent B reads it and queues a reply, and
station A receives that reply despite an interrupted contact and a lost
receipt. A continuous interactive RF session is unnecessary. Global operation
extends this exchange to delayed contact windows and optional relay custody;
it does not assume simultaneous end-to-end connectivity.

The model views below use PlantUML as their version-controlled source. Views
identified as UML describe software structure, interaction, deployment,
activity, and state. The SysML-style block definition (bdd), internal block
(ibd), and requirements views use explicit stereotypes and trace relationships
for system engineering. They are lightweight architectural views, not a
machine-validated SysML model, formal port typing, or proof of standards
conformance. Requirements are proposal obligations, not claims of satisfaction.

\landscapeplantuml{build/diagrams/Linkwise_Chat_BDD}{SysML-style block definition: independently administered stations and their agent, messaging, execution, storage, modem, and RF responsibilities. All messaging assembly is proposed.}{fig:chat-bdd}
\portraitplantuml{build/diagrams/Linkwise_Chat_IBD}{SysML-style internal block view: logical station interfaces from MCP to durable messages, local execution, continuous audio, PTT, and the HF peer. Connectors name contracts rather than formally typed SysML ports.}{fig:chat-ibd}

\subsection{Requirements and traceable evidence}
\label{sec:chat-requirements}
The following requirements define the first service boundary. Concrete limits
are installed station configuration, versioned with the tests and operating
profile; an agent cannot enlarge them by including values in a message.

\begin{longtable}{@{}p{0.13\textwidth}p{0.44\textwidth}p{0.36\textwidth}@{}}
\caption{Proposed chat requirements and acceptance evidence.}\label{tab:chat-requirements}\\
\toprule
\textbf{ID} & \textbf{Required behavior} & \textbf{Acceptance evidence}\\
\midrule
\endhead
CHAT-01 & Return local acceptance only after an atomic durable outbox commit,
including idempotency mapping and initial event. & Crash at each commit boundary;
an accepted message and its identity survive restart.\\
CHAT-02 & Publish final remote receipt only after full validation and an atomic
destination inbox, deduplication, and receipt commit. & Lost-receipt retry yields
one inbox record; conflicting bytes under the same identity are rejected.\\
CHAT-03 & Enforce expiry, storage, body-size, retry, hop, and airtime bounds
independently at every station. & Expiry and exhausted-budget tests stop further
work; cancellation never claims remote deletion.\\
CHAT-04 & Separate claimed address, authenticated principal, station authority,
and received content. & Spoofed bindings and unsupported identity evidence fail
closed; message content cannot authorize CAT, PTT, or tool execution.\\
CHAT-05 & Expose durable, ordered, bounded observations with explicit gaps and
unknown delivery evidence. & Reconnect resumes from cursor or receives an
explicit resynchronization response; stale evidence is not reported as current.\\
CHAT-06 & Reassemble bounded messages and reject malformed or inconsistent
fragments without unbounded allocation. & Reordered, repeated, missing,
oversized, and conflicting fragments have deterministic outcomes.\\
CHAT-07 & Retain local ownership of audio continuity, ACK deadlines, PTT release,
and protocol progress. & Delayed or disconnected MCP does not violate timing,
renew authority, or duplicate contact starts.\\
CHAT-08 & Make relaying optional and distinguish custody from final delivery. &
Multi-hop interruption, loop, capacity, and custody-expiry tests preserve
explicit message location and receipt uncertainty.\\
\bottomrule
\end{longtable}

Figure~\ref{fig:chat-requirements} illustrates the core trace relationships;
Table~\ref{tab:chat-qualification} extends verification to fragmentation,
local execution, and relaying. In the requirements view, a test case's
\texttt{verify} relationship means that it is intended to establish evidence,
not that the test has already passed. The complete requirement registry must
retain all eight identifiers when the implementation design is decomposed.

\landscapeplantuml{build/diagrams/Linkwise_Chat_Requirements}{SysML-style requirements view: selected derived obligations, proposed responsible blocks, and planned verification cases. Satisfaction and verification edges express design allocation, not achieved qualification.}{fig:chat-requirements}

\subsection{Software responsibilities and deployment}
The inbound messaging MCP adapter validates identity, schema, quotas, and
request bounds. A message application service owns acceptance, inbox access,
receipt interpretation, and transactional changes. A durable spool stores
bounded bodies and immutable envelope metadata, while a journal provides
ordered observations. A directory resolves agent addresses to station
bindings with provenance and expiry; route evaluation returns evidence and
exclusions rather than an unsupported assertion that a peer is reachable.

The agent selects objectives and eligible routes or delegates a bounded
selection rule. A contact scheduler executes that installed rule and waits
for eligible windows. It does not become a competing strategic agent.
Station execution retains the existing preparation, lease, configuration
revision, automatic policy-grant, and bounded-start contract. A chat send is
permission to queue a message within the authenticated principal's scope;
it is not an unconditional PTT permit. Scheduling intent and transmission
authority remain separate records.

The modem backend exchanges bounded bytes below the message service. It
must not infer application delivery from a link ACK. A supported Winlink
client adapter is an optional alternative route, with distinct addressing,
service semantics, and confirmation evidence; a generic agent envelope is
not assumed to be accepted or delivered by Winlink without that integration.

Each deployed station has its own radio, antenna, durable state, policy,
identity, agent-host relationship, and administrative failure domain. Several
agents may share one station subject to quotas and fair scheduling. A station
may also receive and retain messages while no agent process is running.
An RF packet does not automatically wake ChatGPT or Codex: a separately
qualified host runner, supported event integration, or polling loop must
arrange agent execution and retain its own processing checkpoints.

\landscapeplantuml{build/diagrams/Linkwise_Chat_Components}{UML component view: the proposed messaging service and contact scheduler sit above the station execution contract; optional Winlink service evidence remains separate from modem-link evidence.}{fig:chat-components}
\landscapeplantuml{build/diagrams/Linkwise_Chat_Deployment}{UML deployment view: illustrative stations across regions, with direct and delayed relay contact possibilities. Edges are candidate opportunities, not validated propagation paths or a guaranteed worldwide mesh.}{fig:chat-deployment}

\subsection{Addresses, envelope, fragmentation, and receipts}
An agent address names an authority namespace and stable agent identifier.
A station identity names an operational endpoint; a directory binding maps
an agent to one or more authorized endpoints at a stated revision and expiry.
A callsign or decoded setup field alone is not application authentication.
The selected identity profile must define provisioning, issuer trust,
binding validation, revocation, and what unauthenticated observations mean.
Cryptographic authentication, payload confidentiality, and RF protocol
compatibility are separate capabilities. Any selected profile must fit the
station's configured service constraints; this chapter does not prescribe an
unqualified over-the-air security scheme.

The immutable end-to-end envelope contains a protocol version, stable message
ID, sender and recipient addresses, conversation ID, optional reply-to message
ID, creation time, expiry time, content type, payload length, and payload
digest. Encoding must define canonical byte order, maximum field lengths,
version negotiation, and the exact bytes covered by identity evidence before
implementation. A digest detects content corruption; it does not establish
sender authenticity. Local filesystem paths, executable instructions, and
unchecked tool commands are not envelope capabilities.

Messages are compact UTF-8 text or explicitly versioned structured payloads.
The body is fragmented below the application boundary with message identity,
byte offset, total bounded length, and fragment consistency checks. The
receiver validates bounds before allocating or writing, permits identical
retransmitted fragments, rejects conflicting overlaps, and publishes no inbox
message until reassembly, digest, expiry, recipient, and identity checks pass.
Persisted range checkpoints may support resumption only where a negotiated
application transport defines them; native modem support does not imply that
such a resume protocol already exists. Link frames remain owned by the modem
session, with a separate application reassembly lifetime and memory budget.

The deduplication key is the validated sender namespace plus message ID.
The associated digest must be immutable: the same key and same digest replays
the existing receipt; the same key and different digest is a conflict.
Deduplication records and tombstones must survive at least the supported
message lifetime plus the configured clock-skew and replay-retention window.
Storage pressure must reject new acceptance or advertise reduced capacity,
not silently discard records still required by that contract. Unknown clocks
or expiry uncertainty require explicit policy decisions and observations.

Receipts bind message ID, digest, issuer, kind, and relevant authentication
and observation evidence. Distinct kinds include local acceptance, relay
custody, final destination inbox receipt, and optional agent processing
acknowledgment. An application's processing receipt may report a result or
rejection; it cannot prove every external side effect occurred exactly once.
Receipt forwarding, retention, retry, and expiry have their own bounded
budget. A transport ACK is never promoted to a stronger receipt kind.

\landscapeplantuml{build/diagrams/Linkwise_Chat_Message_Model}{UML class view: immutable envelopes are distinct from station copies, contact attempts, route evidence, and receipts. Attribute lists describe draft domain data, not a finalized wire serialization.}{fig:chat-message-model}

\subsection{Proposed agent-facing MCP contract}
The names below are new proposal names and are not currently registered MCP
tools. Each mutation carries authenticated principal, request/correlation ID,
idempotency key, policy revision, and a bounded deadline. Retried acceptance
with the same key and same content returns the original result; a changed
request under the same key is rejected. Tool results contain explicit unknown,
stale, unsupported, and terminal reasons rather than inferred success.

\begin{longtable}{@{}p{0.35\textwidth}p{0.58\textwidth}@{}}
\caption{Draft messaging MCP application surface.}\label{tab:chat-tools}\\
\toprule
\textbf{Proposed tool} & \textbf{Contract}\\
\midrule
\endhead
\path{chat_send} & Accept recipient, bounded content or authorized body
reference, conversation/reply references, expiry and eligible delivery intent;
return stable message ID and durable local acceptance. No RF success implied.\\
\path{chat_list_messages} & Return a bounded page of authorized inbox metadata
and opaque body references, with explicit cursor retention and gap behavior.\\
\path{chat_read_message} & Read an authorized bounded body or chunk. Mark content
as received, untrusted data; reading is not processing acknowledgment.\\
\path{chat_get_delivery_status} & Return local work state plus separately
attributed contact, custody, remote-receipt, and processing evidence.\\
\path{chat_get_events} & Return ordered, bounded journal events after a cursor;
on retention gaps require snapshot resynchronization rather than silent loss.\\
\path{chat_cancel} & Stop eligible local work for a message; return cancellation
operation and residual uncertainty. Cannot retract remote or relay copies.\\
\path{chat_ack_processed} & Record the local agent's explicit processing
checkpoint or outcome under its authenticated identity; repeat idempotently.\\
\bottomrule
\end{longtable}

Message bodies stay in a bounded spool or application stream by default.
MCP snapshots and event feeds carry metadata and body references; authorized
body retrieval is explicit and capped. Existing modem observation tools remain
observational. Message intent reaches station preparation and start through
application services, never by exposing raw modem bytes or ad hoc PTT commands
inside \path{chat_send}. An accepted long-running operation returns promptly;
no MCP request must stay open throughout an HF contact. Event cursors are
service-issued opaque tokens, scoped to principal and journal epoch.

Agent hosts retain durable processing checkpoints and reconcile external actions
before retrying them. Replies require turn limits, maximum reply depth,
conversation expiry, per-peer quotas, and an explicit stop outcome. Received
content cannot grant tool access or change station limits.

\portraitplantuml{build/diagrams/Linkwise_Chat_Direct}{UML sequence: durable local acceptance, a bounded direct contact, committed remote inbox receipt, agent polling, and an independently queued reply. Link ACKs and agent processing are separate boundaries.}{fig:chat-direct}

\subsection{State, interruption, cancellation, and restart}
Delivery work and delivery knowledge must not share one ambiguous success
flag. Local work moves through queued, transferring, waiting for receipt,
expired, cancelled, or failed states. Valid remote-receipt and agent-processing
evidence are retained separately and can arrive late. The sender-view state
machine summarizes common progress; a terminal local state never deletes or
suppresses later evidence, nor does such evidence restart transmission.

Cancellation is best-effort stopping of local execution. A contact interrupted
after the last byte may already have produced a remote inbox copy. Expiry stops
future delivery work but cannot undo an earlier receipt. A timeout with no
receipt means uncertainty, not proof of non-delivery. A successful contact
also need not contain a successful application message transfer.

Inbox acceptance, deduplication record, and receipt must be committed in one
transaction or an equivalent recoverable write-ahead protocol. A crash after
commit but before receipt transmission is recovered by replaying the receipt
when the sender retries. A crash before commit cannot produce an authoritative
remote-receipt. Process restart invalidates old preparation handles, leases,
and grants; recovery reconciles device/PTT state and re-evaluates unresolved
messages before any new RF attempt. MCP disconnection alone follows the
already installed operation policy and never extends a permit.

\landscapeplantuml{build/diagrams/Linkwise_Chat_Message_States}{UML sender-view message state machine. Local work termination and remote delivery evidence remain distinct; late receipts are journaled even when no further sending is permitted.}{fig:chat-message-states}
\portraitplantuml{build/diagrams/Linkwise_Chat_Contact_States}{UML lifecycle of one local contact operation, separate from message delivery: fresh authority, establishment, exchange, release, and fault reconciliation. Initial/final nodes delimit the operation; the station owner returns to idle only after release or reconciliation.}{fig:chat-contact-states}
\portraitplantuml{build/diagrams/Linkwise_Chat_Recovery}{UML recovery sequence: restart and a lost receipt cause a policy-checked retry with the same message identity, not duplicate inbox acceptance or an exactly-once agent-action claim.}{fig:chat-recovery}

\subsection{Global contact planning and optional relay custody}
The first route is direct delivery between two configured peers. Worldwide
scope requires asynchronous availability: frequency eligibility, time of day,
antenna/radio capabilities, measured noise, prior contacts, propagation
evidence, and peer listening schedules affect candidate contacts. Propagation
predictions inform ranking; they do not prove reachability or authorize RF.
Each candidate has evidence age, confidence, exclusions, and an expiry. A
scheduler must revalidate these at execution time, honor half-duplex station
ownership, avoid channel contention under its installed access profile, and
apply fairness across agents and conversations.

If no route is eligible, the message remains queued with a next-review time
and an explicit reason, until cancellation, expiry, or a terminal budget
condition. It does not spin in an unbounded autonomous retry loop. Listening
schedules and directory updates also require a bootstrap mechanism: initial
peer addresses, compatible application protocol versions, listening windows,
and relay permissions must be provisioned or learned from a separately
qualified discovery service. No worldwide discovery infrastructure is assumed.

Relaying is an optional later capability, not a consequence of direct modem
interoperability. A relay accepts custody only after validating recipient
routing, identity profile, remaining lifetime, available durable capacity,
per-peer quotas, permitted forwarding scope, and hop history. It records its
commitment and retention/expiry contract before issuing a custody receipt.
The original end-to-end envelope is preserved; bounded, separately validated
hop metadata carries relay identity, visited stations, remaining hop budget,
and the next-hop decision. Repeated station visits, exhausted budgets, and
conflicting identities are rejected. A sender may retain its copy until a
final receipt or until an explicitly chosen custody policy permits release;
there is no deletion merely because a link ACK was received.

Final receipts may need a later reverse contact or another eligible route.
They are deduplicated and rate-limited so acknowledgment traffic cannot become
an unbounded relay loop. Relay custody does not prove delivery to the final
agent or destination station, and a relay failure may leave the sender's
knowledge uncertain. Direct and relay modes must use the same evidence labels
while recording different receipt issuers and copy locations.

\portraitplantuml{build/diagrams/Linkwise_Chat_Routing}{UML activity: a message is sent only after route, contact-window, station, and automatic policy checks; missing evidence leads to bounded deferral rather than presumed global connectivity.}{fig:chat-routing}
\portraitplantuml{build/diagrams/Linkwise_Chat_Relay}{UML relay sequence: durable custody, a later destination contact, and a separately routed final receipt. No continuously connected global path is required or promised.}{fig:chat-relay}

\subsection{Connection to the native DSP and timing pipeline}
The outbound data path is envelope and body, durable spool, application
fragmentation, bounded session byte stream, modem framing/FEC, waveform
synthesis, audio playback, and the radio. The inbound path reverses it through
continuous capture, acquisition, soft decisions, FEC and CRC, link/session
recovery, fragment reassembly, and committed inbox acceptance. The station
interface view in Figure~\ref{fig:chat-ibd} locates the stream and audio
boundaries; the detailed DSP diagrams in Section~\ref{sec:modem-software}
describe the currently implemented algorithms.

There are three distinct integrity scopes: modem CRC validates a frame,
application digest validates a complete reassembled payload, and the selected
identity profile validates whatever sender/issuer assertions it supports.
None substitutes for durable remote receipt. An agent does not choose
symbol boundaries, write audio blocks, or supply ACK timing. Local session
ownership enforces these deadlines independently of MCP latency, model
availability, and network interruptions to the agent host.

Receive overflow, lost audio continuity, playback failure, and unconfirmed
PTT release terminate or fault the affected contact under local policy. The
message service retains unresolved delivery work and evidence rather than
silently stitching samples or treating attempted playback as delivery. The
current CPAL input adapter and synthetic demo are foundations for this path;
output-device integration and qualified operational session assembly remain
prerequisites for an RF chat claim.

\subsection{Qualification gates and open design decisions}
\label{sec:chat-qualification}
Implement direct messaging before routing, and qualify failure handling before
unattended agent replies. Every gate records software/configuration versions,
peer identity, exact payloads, timestamps, receipt provenance, terminal causes,
and local radio/authority observations appropriate to that stage.

\begin{longtable}{@{}p{0.16\textwidth}p{0.38\textwidth}p{0.39\textwidth}@{}}
\caption{Planned messaging qualification gates; none is claimed complete by this draft.}\label{tab:chat-qualification}\\
\toprule
\textbf{Gate} & \textbf{Scope and requirements} & \textbf{Required acceptance evidence}\\
\midrule
\endhead
Q1: durable service & Simulated byte transport and two agent-host runners;
CHAT-01, 02, 04--06. & Exact bidirectional bodies; one committed inbox entry under
retries; crash-point recovery; content conflicts rejected; cursor gaps explicit;
idempotent agent processing and reply-loop bounds exercised.\\
Q2: audio transport & Native synthesized and controlled physical audio path;
CHAT-02, 03, 06, 07. & Audio gaps and corruption cannot become delivery; exact
bodies through full session assembly; lost ACK/receipt, overflow, and playback
faults preserve bounded recovery. Physical audio and synthetic evidence are labeled
separately; any stock-peer qualification is a separate recorded result.\\
Q3: two-station HF & Qualified radio/audio/PTT chain and configured peers;
CHAT-01--07. & A message and reply delivered over RF with traceable final receipts;
interrupted contacts, unavailable peer, station restart, and budget expiry produce
correct local state and release evidence.\\
Q4: relay network & Three or more stations, non-overlapping contact windows;
CHAT-01--08. & Durable custody and later final receipt; loop prevention, relay
capacity rejection, expiry, relay restart, and return-receipt failure tested.\\
\bottomrule
\end{longtable}

Correctness criteria are exact payload bytes, bounded resources, no duplicate
committed inbox entry within the documented deduplication lifetime, and no
promotion of weaker evidence into stronger delivery claims. Report latency
as queue wait, contact establishment, transfer, receipt return, and agent
processing time separately. Measure useful payload goodput, retries,
application overhead, storage occupancy, and keyed airtime. Numerical service
objectives depend on an explicitly qualified station/channel profile; the
whitepaper does not promise a worldwide latency or throughput figure.

Before implementation, resolve the wire encoding and negotiation profile,
identity/address bootstrap and revocation, spool transaction and retention
scheme, clock-skew policy, agent-host wake/poll integration, contact access
profile, and any relay-custody semantics. A versioned interoperability test
suite must fix these choices. The first implementation deliverable is the
small two-agent exchange defined in Section~\ref{sec:chat-scope}, with honest
receipt evidence and failure recovery, rather than an unqualified global mesh.
